% Ingested from kunci_jawaban (tugas_topik_3.tex)
% shared
\begin{aligned}
9 - x \geq 0 \\
-x \geq -9 \\
x \leq 9 \\
D_f = \{x \in \mathbb{R} \mid x \leq 9\} = (-\infty, 9] \\
\sqrt{9 - x} \geq 0 \\
f(9) = \sqrt{9 - 9} = 0
\end{aligned}

% role: concept
\begin{aligned}
f(x) = \sqrt{9 - x} memuat akar kuadrat, sehingga radikan (bilangan di bawah akar) harus non-negatif. \\
radikan non-negatif karena akar kuadrat
\end{aligned}

% role: domain
\begin{aligned}
9 - x \geq 0 \\
-x \geq -9 \\
x \leq 9 \\
D_f = \{x \in \mathbb{R} \mid x \leq 9\} = (-\infty, 9]
\end{aligned}

% role: range
\begin{aligned}
\sqrt{9 - x} \geq 0 \\
f(9) = \sqrt{9 - 9} = 0 \\
Ketika x semakin kecil, 9 - x membesar tanpa batas sehingga f(x) juga membesar tanpa batas. \\
R_f = \{y \in \mathbb{R} \mid y \geq 0\} = [0, \infty).
\end{aligned}

% role: value_table
\begin{aligned}
(-7, 4) \\
(0, 3) \\
(5, 2) \\
(8, 1) \\
(9, 0)
\end{aligned}

% role: figure

% final answer
D_f = (-\infty, 9] ; R_f = [0, \infty)
